\documentclass[12pt,a4paper]{article}

\usepackage[T2A]{fontenc}
\usepackage[utf8]{inputenc}
\usepackage[russian]{babel}
\usepackage{geometry}
\usepackage{setspace}
\usepackage{enumitem}
\usepackage{hyperref}
\usepackage{verbatim}
\usepackage{parskip}

\geometry{left=25mm,right=25mm,top=20mm,bottom=20mm}
\onehalfspacing
\setlist[itemize]{leftmargin=1.25cm}
\setlist[enumerate]{leftmargin=1.25cm}

\begin{document}

\begin{center}
    \textbf{\Large Архитектура продукта}\\[0.5em]
    \textbf{\large Компьютерная игра ``The Hero''}
\end{center}

Ниже приведена архитектура продукта, выведенная только из требований двух приложенных документов: технического задания и описания проекта. Основа продукта в документах формулируется как однопользовательская компьютерная игра для ПК в жанре пошаговой стратегии/тактики с единичной глобальной картой, упрощённым боем без тактической сетки, базой для найма и улучшений, сохранением прогресса, базовыми настройками, локализацией UI и обучением/подсказками. Внешние сервисы, мультиплеер, облачные сохранения и активные AI-соперники в текущей версии не требуются.

\section*{1. Анализ требований из приложенных файлов}

Архитектурно значимое ядро продукта задаётся игровым циклом <<карта $\rightarrow$ событие/бой $\rightarrow$ награда $\rightarrow$ усиление армии/базы>>. Пользователь перемещает героя по клеточной карте, взаимодействует с объектами, вступает в бой с нейтральными противниками, получает ресурсы, опыт и/или артефакты, затем усиливает армию через базу и постройки. Это означает, что архитектура должна строиться не вокруг отдельных экранов, а вокруг единого состояния прохождения, которое последовательно проходит через несколько режимов: глобальная карта, бой, база/найм, информационные окна и настройки.

Из ТЗ следует обязательное наличие следующих функциональных блоков: управление сессией, сохранение прогресса, карта мира, герой и армия, взаимодействие с объектами карты, ресурсы и награды, боевая система, юниты и характеристики, база и найм, интерфейс и настройки, а также локализация и раздел помощи/обучения. Это не позволяет ограничиться только игровым ядром: в архитектуру обязательно должны входить подсистемы сохранений, пользовательских настроек, локализации и помощи.

Карта мира описана как клеточная тайловая сетка, по которой герой перемещается автоматически по кратчайшему доступному маршруту после выбора цели. На карте должны присутствовать охраняемые ресурсы, нейтральные противники, статичные противники-отряды, база и нейтральные строения. Следовательно, архитектуре необходимы: модель карты, модель объектов карты, модуль маршрутизации, модуль обработки достижения цели и механизм фиксации состояний объектов карты.

Боевая система, согласно ТЗ, принципиально не использует тактическую сетку и позиционирование. Бой идёт на отдельном экране, стороны формируются из армии игрока и отряда противника, порядок действий определяется инициативой, игрок выбирает цели для своих отрядов, урон считается без учёта позиции, потери фиксируются, после боя выжившие сохраняются без автоматического восстановления. Отсюда следует, что боевой модуль должен быть независим от карты и работать на собственной модели состояния боя, а данные в него должны передаваться из общего состояния игры в виде снимка состава сторон.

Для базы и найма требования также достаточно конкретны: реализуются 1--3 постройки, каждая отвечает за найм определённого типа юнитов; найм доступен каждые 7 ходов; незабранный найм накапливается; у каждой постройки есть минимум одно улучшение; ранее нанятые юниты базового типа могут быть улучшены за ресурсы до уровня текущей постройки. Значит, архитектуре требуется отдельный модуль недельного цикла, модуль доступного найма, модуль состояния построек и модуль апгрейда юнитов.

Сохранения и надёжность являются отдельным архитектурным требованием, а не вспомогательной функцией. ТЗ требует ручное сохранение, автосохранение не реже одного раза за завершённый игровой день/ход на глобальной карте, резервную копию предыдущего корректного сохранения, валидацию структуры и диапазонов при загрузке, контроль обязательных полей, совместимость в пределах мажорной версии и безопасный откат к последнему корректному состоянию при ошибке. Это означает, что подсистема сохранений должна быть самостоятельной, версионируемой и защищённой от частичных повреждений файла.

Интерфейс тоже задаёт конкретные ограничения: требуются главное меню, экран карты, экран боя, экран базы/найма, окно настроек, информационные окна, управление мышью и частично клавиатурой, поддержка русского и английского UI, настройка звука, а также обучающий режим или подсказки для первых уровней/боёв и настройка сложности. Следовательно, UI нельзя рассматривать как тонкую оболочку: он должен иметь собственную структуру экранов, слой локализации и слой пользовательских настроек.

Из описания проекта дополнительно следует зафиксированный технологический стек: Unity LTS, C\#, uGUI, TextMeshPro, JSON + Newtonsoft.Json, Git/GitHub. Следовательно, речь идёт об архитектуре одного клиентского приложения внутри Unity, а не о распределённой системе или сервисной архитектуре.

\section*{2. Архитектурный подход}

Для данного продукта обоснован модульный монолит внутри Unity. Такое решение непосредственно следует из требований: игра однопользовательская, работает локально на ПК, не использует внешние сервисы, должна хранить прогресс локально и включает ограниченный, но связанный набор игровых режимов. Разделение на отдельные сетевые сервисы, серверную часть или внешние хранилища противоречило бы объёму и формату требований.

Базовый принцип архитектуры должен быть таким: единое доменное состояние игры находится в центре, а режимы карты, боя, базы, настроек и сохранения работают как согласованные подсистемы над этим состоянием. Это необходимо потому, что ТЗ требует сохранять и восстанавливать не только ресурсы и героя, но и состояние карты, армии, прогресса по неделям, побеждённых противников, посещённых объектов и результаты развития базы.

Внутри такого монолита целесообразно разделение на четыре слоя. Первый слой --- доменная модель игры: сущности героя, армии, юнитов, ресурсов, карты, базы, найма, наград и состояния прохождения. Второй слой --- игровые подсистемы: карта, бой, экономика, найм, прогресс, сохранение. Третий слой --- слой представления: экраны и окна UI. Четвёртый слой --- инфраструктурный слой: сериализация JSON, загрузка конфигурации, локализация, аудио, файловый ввод-вывод. Такое разделение нужно, чтобы требования к логике, UI и сохранениям не были смешаны в одних и тех же компонентах.

Отдельно важно, что бой должен быть архитектурно отделён от карты. Документы прямо запрещают моделировать тактическое поле боя, перемещение по клеткам в бою и позиционирование. Поэтому боевой модуль должен получать только состав сторон, инициативу, цели, параметры юнитов и возвращать результат боя, потери и награды. Любое включение пространственной боевой логики было бы отходом от ТЗ.

\section*{3. Компоненты и модули}

\subsection*{3.1. Ядро приложения}

\textbf{Модуль управления приложением и сессией.}

Отвечает за запуск игры, переходы из главного меню, создание новой игры, продолжение сохранённой сессии, завершение работы приложения. Он нужен напрямую из обязательных функций управления сессией.

\textbf{Модуль состояния игры.}

Хранит текущее состояние прохождения как единый агрегат: состояние героя, армии, ресурсов, карты, базы, игровых дней/недель, выбранной сложности, пользовательских настроек, флагов посещения объектов и результатов боёв. Это центральный модуль, потому что практически все требования к сохранению, переходам между режимами и восстановлению прогресса сводятся к одной целостной модели состояния.

\subsection*{3.2. Игровые подсистемы}

\textbf{Модуль карты мира.}

Отвечает за модель тайловой карты, координаты, доступность клеток, размещение объектов карты, отображение точки назначения, расчёт маршрута, расход очков движения и завершение перемещения. Этот модуль обязателен, поскольку в ТЗ отдельно определены клеточная карта, маршрут до цели, точки интереса и трата очков движения.

\textbf{Модуль объектов карты и взаимодействий.}

Обрабатывает достижение выбранного объекта и запускает нужный сценарий: сбор ресурса, переход в бой, выдачу награды, открытие базы, захват или пометку нейтрального строения как использованного/захваченного. Он нужен потому, что сами объекты карты имеют разную природу, но пользователь взаимодействует с ними единым действием --- достижением цели на карте.

\textbf{Модуль героя и армии.}

Хранит параметры героя, состав армии, слоты армии, типы и количество юнитов в каждом отряде, а также ограничения на размер армии. Он необходим как отдельный модуль, потому что армия участвует и в карте, и в базе, и в бою, и в сохранении.

\textbf{Модуль юнитов и характеристик.}

Содержит определения типов юнитов и их параметров: здоровье/численность, атака, защита/броня, урон, инициатива/скорость. Он отделяется от текущего состояния армии, потому что одни и те же типы юнитов используются в разных армиях, в найме, в улучшениях и в расчёте боя.

\textbf{Боевой модуль.}

Включает формирование сторон, очередь ходов по инициативе, приём выбора целей игроком, расчёт урона, учёт потерь, определение победы/поражения и формирование результата боя. Это самостоятельный модуль, потому что бой выполняется в отдельном режиме и должен быть воспроизводим при одинаковых исходных данных.

\textbf{Модуль экономики и ресурсов.}

Учитывает получение ресурсов за карту и бои, расход на найм, улучшения построек и улучшение ранее нанятых юнитов. Он нужен как отдельный модуль, потому что ресурсы участвуют в нескольких сценариях и должны быть единообразно проверяемы при каждом расходовании.

\textbf{Модуль базы, построек и найма.}

Хранит перечень построек базы, их уровень, доступные типы юнитов, накопленный доступный найм, правила пополнения раз в 7 ходов и операции найма/улучшения. Он должен быть выделен отдельно, поскольку механика базы имеет собственный жизненный цикл и собственные данные, не совпадающие ни с картой, ни с боем.

\textbf{Модуль наград и прогресса.}

Формирует и применяет последствия побед и событий: ресурсы, опыт героя, артефакты, изменение характеристик при наличии таких эффектов. Он нужен потому, что ТЗ связывает бой и карту с последующим усилением армии/базы через формализованный шаг <<награда>>.

\textbf{Модуль хода, игрового дня и недели.}

Инкрементирует глобальный ход, завершает игровой день по кнопке <<Завершить ход>>, инициирует автосохранение и пополнение найма по игровой неделе. Этот модуль необходим, так как в документах неделя и автосохранение привязаны именно к завершению глобального хода.

\subsection*{3.3. Подсистемы представления и инфраструктуры}

\textbf{UI-модуль экранов и окон.}

Включает главное меню, экран карты, экран боя, экран базы/найма, окно настроек, окна информации о герое, армии, ресурсах, боях, событиях и постройках. Он нужен, поскольку перечислен в ТЗ как обязательный набор экранов и информационных представлений.

\textbf{Модуль пользовательских настроек.}

Отвечает за громкость, включение/выключение звука, выбор языка интерфейса, уровень сложности, а также настройку скорости анимаций или их отключение там, где это предусмотрено документами. Он требуется напрямую из функциональных и специальных требований.

\textbf{Модуль локализации.}

Обеспечивает русский и английский интерфейс только для UI-элементов. Он должен быть самостоятельным, потому что локализация упомянута как отдельная функция, а стек проекта включает TextMeshPro и uGUI для текстового интерфейса.

\textbf{Модуль помощи/обучения.}

Реализует раздел помощи и обучающий режим или подсказки для первых уровней/первых боёв. Он нужен не как общая дополнительная удобная функция, а потому что прямо указан в требованиях.

\textbf{Модуль сохранения и загрузки.}

Отвечает за сериализацию состояния игры, запись ручных и автоматических сохранений, хранение резервной копии, валидацию структуры при чтении, проверку версии формата и обработку ошибок с безопасным восстановлением. Он является отдельной подсистемой из-за объёма требований к надёжности и совместимости.

\textbf{Модуль конфигурационных данных.}

Загружает параметры юнитов, построек, наград, составов противников и другие конфигурации/баланс из JSON. Он нужен, поскольку ТЗ и описание проекта прямо указывают конфигурационные данные и JSON как формат хранения.

\textbf{Аудиомодуль.}

Обеспечивает проигрывание музыки и звуковых эффектов, а также применение пользовательских настроек громкости и отключения звука. Он требуется, поскольку аудиовыход и управление звуком перечислены в ТЗ.

\section*{4. Связи между модулями}

Управление приложением и UI верхнего уровня инициируют создание новой игры или загрузку сохранения. После этого создаётся или восстанавливается единое состояние игры, которое передаётся в модуль карты как основной режим прохождения.

Модуль карты использует модуль объектов карты и взаимодействий. Когда герой достигает цели, обработчик объекта определяет тип результата: мгновенный сбор ресурса, открытие базы, переход в бой, выдача награды, захват нейтрального строения. Сам модуль карты не должен содержать логику боя или найма; он только инициирует соответствующий сценарий.

При запуске боя модуль взаимодействий передаёт в боевой модуль текущую армию игрока из модуля героя и армии, а также выбранный вражеский отряд с его параметрами из конфигурации/состояния карты. Боевой модуль использует определения юнитов и возвращает результат: потери сторон, исход боя, награду и обновлённое состояние армии. Затем модуль наград и прогресса применяет последствия к общему состоянию игры.

Модуль базы, построек и найма получает доступ к модулю экономики и ресурсов для проверки стоимости, к модулю героя и армии для добавления юнитов в слоты и к модулю хода/недели для расчёта накопленного найма. Таким образом, база не может существовать изолированно: она зависит от времени, ресурсов и текущего состава армии.

Модуль хода и недели воздействует сразу на несколько подсистем: обновляет счётчики времени, инициирует недельное пополнение найма в базе, запускает автосохранение и фиксирует завершение игрового дня. Это один из ключевых координационных модулей всей архитектуры.

Модуль сохранения и загрузки работает не с отдельными экранами, а с единым состоянием игры и пользовательскими настройками. При сохранении он получает снимок текущего состояния. При загрузке он выполняет валидацию, проверку версии и либо восстанавливает корректное состояние, либо предлагает резервное сохранение/новую игру.

UI-модуль связан со всеми игровыми подсистемами только через чтение их состояния и вызов разрешённых действий. Это архитектурно важно: интерфейс должен отображать данные и передавать команды, но не хранить собственную игровую логику, иначе трудно обеспечить консистентное сохранение и корректные переходы между режимами.

\section*{5. Данные и сущности}

Ключевой агрегат данных --- \textbf{состояние игры}. В него должны входить: состояние героя, состав армии, ресурсы, состояние карты, состояние базы, текущий игровой день/неделя, результаты посещения объектов, настройки сложности и данные, необходимые для сохранения/восстановления прогресса. Такой состав следует напрямую из требований к сохранениям и игровому циклу.

\textbf{Карта} должна представляться как набор клеток и размещённых на них объектов. Для объектов нужны по меньшей мере тип, координаты и состояние. Состояние объекта должно позволять различать, например, собранный ресурс, побеждённого противника, посещённый или захваченный объект. Это требование вытекает из разделов о карте, взаимодействиях и сохранении.

\textbf{Герой} --- это сущность с параметрами, участвующая в перемещении по карте и получении наград. Документы явно фиксируют наличие параметров героя, опыта и артефактов, изменяющих характеристики, но не задают исчерпывающую модель навыков; поэтому в архитектуре герой должен рассматриваться как сущность с изменяемыми параметрами и накоплением опыта, без детализации того, чего документы не определяют.

\textbf{Армия} должна храниться как набор слотов, где каждый слот содержит один отряд. \textbf{Отряд} --- это тип юнита и его текущее количество. Такое представление обязательно, потому что ТЗ фиксирует ограничения по слотам и описывает армию именно как набор отрядов одного типа.

\textbf{Тип юнита} должен хранить боевые характеристики: здоровье/численность, атаку, защиту/броню, урон, инициативу/скорость. Отдельно от этого должен храниться \textbf{экземпляр отряда}, содержащий текущую численность, чтобы после боя можно было зафиксировать выживших без автоматического восстановления.

\textbf{База} должна включать перечень построек. \textbf{Постройка} должна содержать тип, уровень, доступный тип юнита для найма, объём доступного найма на текущий момент и данные для улучшения. Требование на накопление найма означает, что у постройки должен быть сохраняемый счётчик накопленного предложения, а не только текущая доступность на экране.

\textbf{Ресурсы} должны храниться в виде отдельной сущности-кошелька, потому что используются в нескольких независимых сценариях: сбор на карте, награда после боя, найм, улучшение построек, улучшение ранее нанятых юнитов.

\textbf{Награда} должна быть отдельной сущностью, объединяющей ресурсы, опыт героя и артефакты. Это нужно потому, что документы прямо допускают награду смешанного состава, а не только однотипную.

\textbf{Сохранение} должно содержать данные о версии формата, обязательных полях, текущем состоянии игры и, желательно, ссылочно независимую структуру, пригодную для проверки целостности и миграции в пределах одной мажорной версии. Это вытекает из требований к валидации и совместимости сохранений.

\textbf{Пользовательские настройки} должны быть отделены от основного игрового состояния и включать по меньшей мере звук, язык интерфейса, сложность и параметры анимации в пределах реально поддерживаемых функций.

\section*{6. Сценарии работы системы}

\subsection*{Сценарий 1. Запуск и вход в сессию}

Пользователь запускает приложение, попадает в главное меню, выбирает новую игру либо продолжение. При новой игре инициализируется начальное состояние карты, героя, армии, базы и ресурсов. При загрузке читается сохранение, проходит валидация, затем состояние восстанавливается. При ошибке загрузки пользователь получает сообщение и вариант отката к резервной копии или начала новой игры.

\subsection*{Сценарий 2. Перемещение по карте}

Пользователь выбирает клетку или объект. Модуль карты рассчитывает кратчайший доступный маршрут, списывает очки движения и перемещает героя до достижения цели или исчерпания движения. По достижении цели управление передаётся модулю взаимодействий.

\subsection*{Сценарий 3. Взаимодействие с ресурсом или объектом карты}

Если достигнут объект-ресурс, ресурс добавляется в экономику, а объект меняет состояние на собранный. Если достигнута база, открывается экран базы/найма. Если достигнут охраняемый объект или противник, запускается бой. Если достигнуто нейтральное строение, выполняется его сценарий взаимодействия и объект помечается как использованный или захваченный в зависимости от реализации документа.

\subsection*{Сценарий 4. Бой}

В боевой режим передаются армия игрока и вражеский отряд. Формируется порядок ходов по инициативе. Игрок выбирает цель для своего отряда, выполняется расчёт урона, фиксируются потери. После завершения боя определяется победа или поражение, выжившие отряды возвращаются в общее состояние игры, победителю начисляется награда. Пространственная модель боя при этом не используется.

\subsection*{Сценарий 5. Усиление армии и базы}

После посещения базы пользователь открывает список построек, видит доступный найм и уровни построек, тратит ресурсы на найм либо улучшение. При наличии улучшенной постройки возможен апгрейд ранее нанятых юнитов соответствующего типа до текущего уровня. Итоговые изменения применяются одновременно к ресурсам, состоянию базы и армии.

\subsection*{Сценарий 6. Завершение хода}

Пользователь завершает игровой ход на глобальной карте. Система фиксирует выполненные действия, обновляет состояние карты, армии и героя, увеличивает счётчик дня, проверяет смену игровой недели, при необходимости пополняет доступный найм и создаёт автосохранение с резервной копией предыдущего состояния.

\subsection*{Сценарий 7. Работа с настройками, локализацией и обучением}

Пользователь открывает настройки, меняет громкость, включает или выключает звук, выбирает язык интерфейса, сложность и, при наличии, параметры анимации. Отдельно доступны раздел помощи/обучения или подсказки для первых боёв/уровней. Эти параметры должны применяться без нарушения основного игрового состояния.

\section*{7. Схема архитектуры}

\begin{verbatim}
[Управление приложением / Главное меню]
                |
                v
      [Единое состояние игры]
                |
    +-----------+-----------+----------------+-----------------+
    |                       |                |                 |
    v                       v                v                 v
[Карта мира]      [База / постройки / найм] [Сохранение]   [Настройки]
    |                       |                |                 |
    v                       v                v                 v
[Объекты карты]     [Экономика / ресурсы] [JSON + backup] [Звук / язык /
    |                       |                                  сложность /
    +-----------+-----------+                                  анимации]
                |
                v
        [Герой и армия]
                |
                v
        [Юниты и характеристики]
                |
                v
           [Боевой модуль]
                |
                v
      [Награды / опыт / артефакты]
                |
                v
      [Обновление общего состояния]
                |
                v
      [UI экранов и информационных окон]
\end{verbatim}

Смысл схемы в том, что карта, бой, база и сохранение не существуют как независимые мини-приложения, а работают поверх одного и того же состояния игры. Именно это позволяет выполнить требование на полный цикл <<карта $\rightarrow$ бой $\rightarrow$ награда $\rightarrow$ усиление армии/базы>> и корректное восстановление прогресса.

\section*{8. Риски и ограничения}

Главный риск --- неполная определённость боевой модели. Документы фиксируют инициативу, выбор цели, расчёт урона и отсутствие позиционирования, но не задают точные формулы урона, правила очередности при равной инициативе, детализацию раундов и степень участия героя в бою. Это влияет на границы боевого модуля и на формат конфигурационных данных.

Второй риск --- неполная определённость модели героя и прогрессии. В одном месте документы говорят об опыте героя и артефактах, в сравнительной части упоминают развитие героя уровней/навыков, но в обязательных функциях нет точной спецификации уровней, навыков, дерева развития и формулы влияния героя на армию. Архитектурно нужно предусмотреть расширяемую сущность героя, но конкретную модель прогрессии документы не задают.

Третий риск --- артефакты описаны неоднозначно. В ТЗ ресурсы в одном фрагменте могут выступать в роли артефактов, а в сравнительных критериях отдельно фигурируют артефакты и инвентарь/экипировка артефактов. При этом как обязательная функция инвентарь артефактов прямо не закреплён. Следовательно, наличие самих артефактов поддержать нужно, а наличие отдельной экипировки/слотов является спорным местом.

Четвёртый риск --- механика нейтральных строений и захвата определена недостаточно. Указано, что такие строения могут быть захвачены и давать бонус/ресурс/доступ к найму/эффекту, но конкретные типы эффектов, длительность действия и повторяемость использования не зафиксированы. Архитектуре нужен обобщённый механизм сценариев объектов, но конкретный набор сценариев требует уточнения.

Пятый риск --- сложность и обучение заявлены, но их точное содержание не определено. Для сложности не сказано, какие именно параметры меняются: сила противников, награды, экономические коэффициенты или иное. Для обучения не определён формат: отдельный режим, встроенные подсказки или сценарий первых боёв. Поэтому соответствующие модули должны проектироваться как конфигурируемые, а не жёстко прошитые.

Шестой риск --- не полностью определён формат журнала событий. В документах есть термин <<журнал событий>> и указание на сообщения/результаты боёв и событий, но не зафиксировано, нужен ли именно отдельный постоянный экран журнала или достаточно информационных сообщений и окон результатов. Это нужно уточнить до детализации UI-архитектуры.

Седьмой риск --- временной цикл нуждается в уточнении. Зафиксировано, что игровой день завершается кнопкой <<Завершить ход>>, неделя состоит из 7 дней, а найм пополняется по неделям. Но не определено, как именно обновляются очки движения героя и какие ещё игровые параметры меняются на границе дня. Для архитектуры это критично, потому что влияет на модуль глобального хода.

Восьмое ограничение --- продукт должен оставаться локальным клиентским приложением без внешних интеграций. Это означает, что любые облачные, сетевые или многопользовательские расширения не должны закладываться в базовую архитектуру как обязательные компоненты текущей версии.

\section*{9. Итоговый вывод}

Архитектура данного продукта должна быть построена как модульный монолит Unity-приложения с единым состоянием игры в центре и с чётким разделением на подсистемы карты, взаимодействий, боя, армии, ресурсов, базы/найма, хода/недели, сохранений, UI, настроек, локализации и обучения. Такой состав компонентов прямо вытекает из ТЗ и описания проекта и соответствует именно вашему продукту: одной карте, нейтральным противникам, бою без тактической сетки, накопительному найму по неделям, локальным сохранениям с резервированием и базовым настройкам.

Ключевая архитектурная особенность проекта --- не жанр вообще, а связка трёх режимов через единое состояние прохождения: глобальная карта, отдельный боевой экран и экран базы/найма. Именно поэтому центральными в архитектуре являются не сцены Unity сами по себе, а модель игрового состояния, переходы между режимами и надёжная подсистема сохранения/восстановления.

Наиболее важные уточнения перед дальнейшей детализацией: точная формула боя, модель прогрессии героя, статус артефактов и инвентаря, конкретика нейтральных строений, механика сложности, формат обучения и состав событийного журнала. До их уточнения общая архитектурная схема уже определима, но часть внутренних правил модулей должна считаться предварительной.

\end{document}